\documentclass[10pt,a4paper]{amsart}
\usepackage{amsmath,amssymb,amsthm,graphicx}
\usepackage{verbatim}
\newenvironment{reference}{\comment}{\endcomment}
%\newenvironment{reference}{}{}
\newenvironment{slogan}{\comment}{\endcomment}
\newenvironment{history}{\comment}{\endcomment}

% For commutative diagrams we use Xy-pic
\usepackage[all]{xy}

% We use 2cell for 2-commutative diagrams.
\xyoption{2cell}
% Theorem environments.
%
\theoremstyle{plain}
\newtheorem{theorem}[subsection]{Theorem}
\newtheorem{proposition}[subsection]{Proposition}
\newtheorem{lemma}[subsection]{Lemma}

\theoremstyle{definition}
\newtheorem{definition}[subsection]{Definition}
\newtheorem{example}[subsection]{Example}
\newtheorem{exercise}[subsection]{Exercise}
\newtheorem{situation}[subsection]{Situation}

\theoremstyle{remark}
\newtheorem{remark}[subsection]{Remark}
\newtheorem{remarks}[subsection]{Remarks}

\numberwithin{equation}{subsection}

% Macros
%
\def\lim{\mathop{\mathrm{lim}}\nolimits}
\def\colim{\mathop{\mathrm{colim}}\nolimits}
\def\Spec{\mathop{\mathrm{Spec}}}
\def\Hom{\mathop{\mathrm{Hom}}\nolimits}
\def\Ext{\mathop{\mathrm{Ext}}\nolimits}
\def\SheafHom{\mathop{\mathcal{H}\!\mathit{om}}\nolimits}
\def\SheafExt{\mathop{\mathcal{E}\!\mathit{xt}}\nolimits}
\def\Sch{\mathit{Sch}}
\def\Mor{\mathop{\mathrm{Mor}}\nolimits}
\def\Ob{\mathop{\mathrm{Ob}}\nolimits}
\def\Sh{\mathop{\mathit{Sh}}\nolimits}
\def\NL{\mathop{N\!L}\nolimits}
\def\CH{\mathop{\mathrm{CH}}\nolimits}
\def\proetale{{pro\text{-}\acute{e}tale}}
\def\etale{{\acute{e}tale}}
\def\QCoh{\mathit{QCoh}}
\def\Ker{\mathop{\mathrm{Ker}}}
\def\Im{\mathop{\mathrm{Im}}}
\def\Coker{\mathop{\mathrm{Coker}}}
\def\Coim{\mathop{\mathrm{Coim}}}

% Boxtimes
%
\DeclareMathSymbol{\boxtimes}{\mathbin}{AMSa}{"02}

%
% Macros for moduli stacks/spaces
%
\def\QCohstack{\mathcal{QC}\!\mathit{oh}}
\def\Cohstack{\mathcal{C}\!\mathit{oh}}
\def\Spacesstack{\mathcal{S}\!\mathit{paces}}
\def\Quotfunctor{\mathrm{Quot}}
\def\Hilbfunctor{\mathrm{Hilb}}
\def\Curvesstack{\mathcal{C}\!\mathit{urves}}
\def\Polarizedstack{\mathcal{P}\!\mathit{olarized}}
\def\Complexesstack{\mathcal{C}\!\mathit{omplexes}}
% \Pic is the operator that assigns to X its picard group, usage \Pic(X)
% \Picardstack_{X/B} denotes the Picard stack of X over B
% \Picardfunctor_{X/B} denotes the Picard functor of X over B
\def\Pic{\mathop{\mathrm{Pic}}\nolimits}
\def\Picardstack{\mathcal{P}\!\mathit{ic}}
\def\Picardfunctor{\mathrm{Pic}}
\def\Deformationcategory{\mathcal{D}\!\mathit{ef}}
\usepackage{cleveref}
\crefname{originalsection}{original section}{original sections}
\Crefname{originalsection}{Original section}{Original sections}
\makeatletter
\newlabel{formal-spaces-section-completion}{{0AIX}{source}{}{sourcepointer-1}{}}
\newlabel{formal-spaces-section-completion@cref}{{[originalsection][0][]0AIX}{[1][0][]source}{}{}{}}
\makeatother
\crefname{originallemma}{original lemma}{original lemmas}
\Crefname{originallemma}{Original lemma}{Original lemmas}
\makeatletter
\newlabel{spaces-more-morphisms-lemma-check-smoothness-on-infinitesimal-nbhds}{{0APQ}{source}{}{sourcepointer-2}{}}
\newlabel{spaces-more-morphisms-lemma-check-smoothness-on-infinitesimal-nbhds@cref}{{[originallemma][0][]0APQ}{[1][0][]source}{}{}{}}
\makeatother
\crefname{originallemma}{original lemma}{original lemmas}
\Crefname{originallemma}{Original lemma}{Original lemmas}
\makeatletter
\newlabel{spaces-morphisms-lemma-etale-universally-injective-open}{{05W5}{source}{}{sourcepointer-3}{}}
\newlabel{spaces-morphisms-lemma-etale-universally-injective-open@cref}{{[originallemma][0][]05W5}{[1][0][]source}{}{}{}}
\makeatother
\title{Factoring a morphism through a monomorphism via the open part and the completion}
\author{The Stacks Project Authors}
\date{stacks-project:restricted.tex:local-snapshot-20261012; GNU FDL 1.2 or later}
\begin{document}
\maketitle
\noindent\textit{Editorial note.} All delivered mathematical text is the author's own, copied from the pinned Stacks Project chapter file restricted.tex; no statement or proof step was rewritten, reordered or removed.\par\medskip
\noindent\textit{Editorial note.} The counted proof is the author's complete argument: the proof builds the fibre product, shows its projection is etale over the closed part via the infinitesimal criterion and is a monomorphism, identifies it as an open immersion, and concludes by the covering of |X|\par\medskip
\noindent\textit{Editorial note.} The label lemma-factor names the counted statement, and the proof environment immediately following it in the source is delivered verbatim as its proof.\par\medskip
\noindent\textit{Editorial note.} The counted statement is carried by the author's own lemma environment in the chapter source.\par\medskip
\noindent\textit{Allowed standard prerequisites (editorial).}\begin{itemize}
\item the completion functor section: morphisms with image in T factor through X\_\{/T\}-related fibre products (Formal Spaces, Section formal-spaces-section-completion)
\item the infinitesimal criterion for etaleness (More on Morphisms of Spaces, Lemma spaces-more-morphisms-lemma-check-smoothness-on-infinitesimal-nbhds)
\item an etale monomorphism of algebraic spaces is an open immersion (Morphisms of Spaces, Lemma spaces-morphisms-lemma-etale-universally-injective-open)
\end{itemize}\medskip
\ifdefined\hypertarget\hypertarget{sourcepointer-1}{}\fi
\noindent\textit{Source reference.} Original section 0AIX: sources/stacks-chow/tags.txt, source lines 13280--13280. Prerequisite pointer invoked at proof line(s) 6472; it names a standard fact of a sibling chapter that the counted argument uses, and every step of the counted proof built on it is retained verbatim in the delivered body.\par\medskip
\ifdefined\hypertarget\hypertarget{sourcepointer-2}{}\fi
\noindent\textit{Source reference.} Original lemma 0APQ: sources/stacks-chow/tags.txt, source lines 13483--13483. Prerequisite pointer invoked at proof line(s) 6474; it names a standard fact of a sibling chapter that the counted argument uses, and every step of the counted proof built on it is retained verbatim in the delivered body.\par\medskip
\ifdefined\hypertarget\hypertarget{sourcepointer-3}{}\fi
\noindent\textit{Source reference.} Original lemma 05W5: sources/stacks-chow/tags.txt, source lines 7574--7574. Prerequisite pointer invoked at proof line(s) 6479; it names a standard fact of a sibling chapter that the counted argument uses, and every step of the counted proof built on it is retained verbatim in the delivered body.\par\medskip
\begin{lemma}
\label{lemma-factor}
Let $S$ be a scheme. Let $f : X \to Y$ and $g : Z \to Y$ be morphisms
of algebraic spaces. Let $T \subset |X|$ be closed.
Assume that
\begin{enumerate}
\item $X$ is locally Noetherian,
\item $g$ is a monomorphism and locally of finite type,
\item $f|_{X \setminus T} : X \setminus T \to Y$ factors through $g$, and
\item $f_{/T} : X_{/T} \to Y$ factors through $g$,
\end{enumerate}
then $f$ factors through $g$.
\end{lemma}

\begin{proof}
Consider the fibre product $E = X \times_Y Z \to X$.
By assumption the open immersion $X \setminus T \to X$
factors through $E$ and any morphism $\varphi : X' \to X$ with
$|\varphi|(|X'|) \subset T$ factors through $E$ as well, see
Formal Spaces, Section \ref{formal-spaces-section-completion}.
By More on Morphisms of Spaces, Lemma
\ref{spaces-more-morphisms-lemma-check-smoothness-on-infinitesimal-nbhds}
this implies that $E \to X$ is \'etale at every point of $E$
mapping to a point of $T$. Hence $E \to X$ is an \'etale
monomorphism, hence an open immersion
(Morphisms of Spaces, Lemma
\ref{spaces-morphisms-lemma-etale-universally-injective-open}).
Then it follows that $E = X$ since our assumptions imply that $|X| = |E|$.
\end{proof}

\end{document}
